\documentclass[12pt, a4paper]{article}

\usepackage[utf8]{inputenc}
\usepackage[T1]{fontenc}
\usepackage[english]{babel}
\usepackage{lmodern}
\usepackage[top=3cm, bottom=2cm, left=3cm, right=2cm]{geometry}
\usepackage{setspace}
\onehalfspacing
\usepackage{indentfirst}
\usepackage[numbers]{natbib}
\usepackage[hidelinks]{hyperref}
\usepackage[hidelinks]{hyperref}
\hypersetup{
    pdftitle={Modular Evaluation and Enhancement of VillageNet for Unsupervised Clustering},
    pdfauthor={Arthur Henrique Moreira Santos}
}
\title{\textbf{Modular Evaluation and Enhancement of VillageNet for Unsupervised Clustering}}
\author{Arthur Henrique Moreira Santos \\ \small UFSJ}
\date{}

\begin{document}

\maketitle

\section{Related Work}

Unsupervised clustering in complex and biomedical datasets has driven the development of graph-based approaches capable of handling high dimensionality and nonlinear 
boundaries without requiring the number of clusters ($k$) to be specified a priori. In this context, the foundational work for this research is \citet{ballal2026villagenet}, 
who introduced the \textbf{VillageNet} algorithm. The authors demonstrated that traditional methods based on $k$-nearest neighbor ($k$-NN) graphs combined with modularity 
optimization, such as the \textbf{Louvain} algorithm \citep{blondel2008louvain}, suffer from distance concentration in high-dimensional spaces and tend to over-fragment 
the data. To overcome this limitation, VillageNet partitions the feature space into locally linear subregions (Voronoi cells) and applies the random-walk-based community 
detector \textbf{WLCF} \citep{ballal2022wlcf} over the reduced inter-village graph, achieving linear time complexity and outperforming classical methods on both general 
and medical benchmarks.

Despite VillageNet's efficiency, its initial partitioning stage relies on the classical \textbf{K-Means} algorithm \citep{macqueen1967some} with a single random 
initialization and Euclidean distance, making it sensitive to local minima and outliers. In the partitional clustering literature, several extensions have been 
proposed to address these issues and serve as the foundation for the variations investigated in this study. \citet{arthur2007kmeanspp} introduced \textbf{K-Means++}, a
 probabilistic seeding method that distributes initial centroids more effectively, reducing variance and improving clustering accuracy. Aiming at scalability for 
 large-scale datasets, \citet{sculley2010minibatch} proposed \textbf{MiniBatch K-Means}, which updates centroids using small random batches, substantially reducing 
 computational cost with minimal loss in quality. For noisy scenarios, \citet{ng2002clarans} consolidated medoid-based algorithms (\textbf{K-Medoids / CLARANS}), which 
 select actual data points as cluster centers, providing greater robustness against extreme values frequently encountered in clinical records.

Regarding the community detection stage on weighted graphs, the literature offers well-established alternatives to WLCF that can be explored on the reduced inter-village 
graph. \citet{traag2019leiden} proposed the \textbf{Leiden} algorithm as a direct refinement of Louvain, resolving the issue of internally disconnected communities while 
ensuring faster convergence and well-connected partitions. Following a premise similar to WLCF---based on random walks---\citet{rosvall2008infomap} developed \textbf{Infomap}, 
which identifies communities by minimizing the description length of information flow across the network via the map equation. Finally, \citet{dhillon2004spectral} 
established the equivalence between \textbf{Spectral Clustering} and normalized graph cuts, a technique widely used to capture non-convex structures through the 
eigenvectors of the Laplacian matrix, whose computational cost becomes tractable when applied to the reduced graph of villages.

Unlike the original study by \citet{ballal2026villagenet}, which validates VillageNet's fixed architecture by comparing it externally against other clustering algorithms, 
this work focuses on the modular evaluation and optimization of VillageNet's internal structure. By systematically integrating and comparing partitional methods 
(K-Means++, MiniBatch K-Means, and K-Medoids), continuous boundary edge weighting, and alternative community detection algorithms (Leiden, Infomap, and Spectral 
Clustering) within the VillageNet pipeline, this research aims to identify the most accurate and stable architectural configuration for both public benchmark 
datasets and real-world clinical trajectories of patients with Acute Myocardial Infarction (AMI).

\bibliographystyle{unsrtnat}
\bibliography{references}

\end{document}